% Generated draft: requires mathematical and source review.
\documentclass[11pt]{article}
\usepackage[margin=25mm]{geometry}
\usepackage{fontspec}
\setmainfont{Latin Modern Roman}
\usepackage{amsmath,amssymb,mathtools,unicode-math}
\setmathfont{Latin Modern Math}
\usepackage{enumitem,needspace,xurl,hyperref}
\setlength{\parindent}{0pt}
\setlength{\parskip}{5pt}
\setlength{\emergencystretch}{3em}
\newcommand{\N}{\mathbb{N}}
\newcommand{\Z}{\mathbb{Z}}
\newcommand{\Q}{\mathbb{Q}}
\newcommand{\R}{\mathbb{R}}
\newcommand{\C}{\mathbb{C}}
\newcommand{\F}{\mathbb{F}}
\newcommand{\A}{\mathbb{A}}
\newcommand{\PP}{\mathbb{P}}
\newcommand{\E}{\mathbb{E}}
\newcommand{\eps}{\varepsilon}
\newcommand{\dd}{\,\mathrm{d}}
\begin{document}
\section*{UnsolvedMath ID: 30003085}
\textbf{Source problem code:} catalogue entry
\subsection{Definitions and mathematical statement}
Let \(\mathbb{P}^{2}(\R)\) be the real projective plane, let \(P\subset\mathbb{P}^{2}(\R)\) be finite, and let \(d\in\Z_{\geq 1}\). Put
\[
V_d=\R[X,Y,Z]_d,\qquad N_d=\binom{d+2}{2}-1=\frac{d^2+3d}{2},
\]
where \(V_d\) is the vector space of homogeneous polynomials of degree \(d\). For \(0\ne F\in V_d\), define
\[
C_F=\{[x:y:z]\in\mathbb{P}^{2}(\R):F(x,y,z)=0\}.
\]
For a finite set \(S\subset\mathbb{P}^{2}(\R)\), define
\[
I_d(S)=\{F\in V_d:F(p)=0\text{ for every }p\in S\}.
\]
Vanishing is independent of the choice of homogeneous coordinates for \(p\).

\textbf{Supplied conjecture, with an explicit conditional interpretation.} If a degree-\(d\) curve means a nonzero degree-\(d\) homogeneous equation up to multiplication by a nonzero real scalar, with reducible and nonreduced equations permitted, the supplied assertion becomes: for every finite \(P\subset\mathbb{P}^{2}(\R)\) and every integer \(d\geq1\), either
\[
I_d(P)\ne\{0\},
\]
or there exist \(Q\subseteq P\) and \(0\ne F\in V_d\) such that
\[
|Q|=N_d,\qquad P\cap C_F=Q,\qquad I_d(Q)=\R F.
\]
Thus the second alternative requires exactly \(N_d\) points of \(P\) on the curve, and those points determine its equation uniquely up to nonzero scalar multiplication.

\textbf{Unresolved qualifications.} The source record does not specify whether curves must be reduced or irreducible, whether nonreduced equations are admitted, or whether uniqueness concerns equations, algebraic curves, or their real point sets. The displayed formulation makes the coefficient-space interpretation precise but requires confirmation against the intended source convention. These interpretations need not coincide. The record gives no lower bound on \(|P|\), no general-position hypothesis, and no exceptions in \(d\); none have been added.
\subsection{Short English statement}
The supplied conjecture asks whether, for every finite set of real projective points and every positive degree \(d\), either all the points lie on one degree-\(d\) curve or there is a degree-\(d\) curve containing exactly \((d^2+3d)/2\) of them and uniquely determined by those points. The meaning of curve and uniqueness needs confirmation; the mathematical formulation above explicitly uses homogeneous equations up to scalar.
\subsection{Sources}
\begin{itemize}
\item Supplied local UnsolvedMath registry record, permanent ID \(30003085\), title: Ordinary Degree-\(d\) Curves Through Finite Point Sets; source problem code \texttt{catalogue entry}. The record belongs to the Oberwolfach Reports Open Problems collection and gives proposed year 2016. No individual proposer is supplied.
\item Supplied primary report citation: \emph{Mini-Workshop: Arrangements of Subvarieties, and their Applications in Algebraic Geometry} (2016), Oberwolfach Reports, pp.~653--702, DOI 10.4171/owr/2016/14. \url{https://doi.org/10.4171/owr/2016/14}. The supplied background describes Oberwolfach Reports as documenting mathematical workshops and their open-problem discussions.
\item Supplied literature source: Boys--Valculescu--De Zeeuw. \url{https://doi.org/10.1016/j.ejc.2022.103509}. No article title or precise theorem statement is supplied in this record.
\item Supplied literature source: Czapliński et al. \url{https://numdam.org/item/10.4171/RSMUP/136-13.pdf}. No article title or precise theorem statement is supplied in this record.
\item Literature assessment dated exactly \texttt{2026-08-21}: the supplied curation labels the problem partially solved and reports that the real-plane conjecture is proved in the conic case, that a cubic analogue is proved for sufficiently large point sets, and that the all-degree statement and sharp finite-set formulation remain open. This is supplied dated curation derived from the OWR record, not an independently verified literature assessment. The supplied background expressly requests independent verification of time-sensitive status and cited claims.
\item The registry labels the statement \texttt{corrected\_verified}, explaining that notation, typography, and sentence structure were repaired without changing mathematical content. This supplied editorial label does not establish mathematical correctness or independent source verification here.
\end{itemize}
\end{document}
